\chapter{坐标环}\label{chap:coordinate-rings}

\noindent\textbf{第 \thechapter{} 章学习目标}
\begin{itemize}
  \item 分别用多项式函数和商环定义仿射簇 $X$ 的坐标环 $K[X]$.
  \item 定义 $K$-代数, 并说明怎样把有限生成的 $K$-代数与多项式环的商等同起来.
  \item 判定一个环是否既约, 并借助商环判定理想是否为根理想.
  \item 用代数性质刻画坐标环: 它们正是有限生成的既约 $K$-代数.
  \item 在具体例子中, 给定一个有限生成的既约 $K$-代数, 求以它为坐标环的仿射簇.
\end{itemize}

第~\ref{chap:varieties-and-ideals}~章的工作给出了在代数与几何两界之间往返的方法, 其中 $\mathcal{V}$ 和 $\mathcal{I}$ 算子是 $\mathbb{A}^{n}$ 中的仿射簇与 $K[x_{1},\ldots,x_{n}]$ 中的根理想之间互逆的双射. 但在代数一侧, 理想还有一个我们尚未用到的特殊作用: 正是对 $K[x_{1},\ldots,x_{n}]$ 的这些子集, 可以取商得到一个环. 若 $X\subseteq\mathbb{A}^{n}$ 是一个仿射簇, 那么应当怎样从仿射簇 $X$ 的角度解释 $K[x_{1},\ldots,x_{n}]$ 对 $\mathcal{I}(X)$ 取商所得的环?

如本章所见, 答案是这个商环自然同构于 $X$ 的坐标环, 即由从 $X$ 到基域 $K$ 的多项式函数所组成的环 $K[X]$. 在第~\ref{sec:polynomial-functions-on-affine-varieties}~节给出这些对象的精确定义后, 我们将有一种从几何走向代数的新方法:
\[
  \begin{aligned}
    \{\text{仿射簇}\}&\to\{\text{环}\}\\
    X&\mapsto K[X].
  \end{aligned}
\]
照例, 我们要问这个对应是否是一本双向词典. 一开始并不是, 因为并非每个环都是某个 $X$ 的 $K[X]$. 为了从代数上刻画哪些环可以成为坐标环, 我们将定义 $K$-代数这一概念, 并研究它的关键代数性质. $K$-代数是一类特殊的环, 包括多项式环及其商环. 本章最后得出的结论是: 一个环恰好在它是有限生成的既约 $K$-代数时才是某个仿射簇的坐标环.

\par\smallskip
{\footnotesize
\noindent\textcopyright{} The Author(s) 2025\par
\noindent E. Clader, D. Ross, \emph{Beginning in Algebraic Geometry}, Undergraduate Texts in Mathematics,
\href{https://doi.org/10.1007/978-3-031-88819-9_4}{\nolinkurl{https://doi.org/10.1007/978-3-031-88819-9_4}}.\par}
